\documentclass[11pt,a4paper]{article}
\usepackage[margin=2.2cm]{geometry}
\usepackage{booktabs,amsmath,hyperref,enumitem,xcolor}
\hypersetup{colorlinks,urlcolor=blue!50!black,linkcolor=blue!50!black}
\setlist{itemsep=1pt,topsep=3pt}
\title{\textbf{CrackDistill: project outline, status and next steps}\\\large Knowledge distillation from SAM models into YOLO for crack segmentation}
\author{Shahin Alakbarov}\date{12 November 2026}
\begin{document}\maketitle

\section{Project outline}
I want a method that transfers what SAM foundation models know into a small YOLO segmenter that runs in real time on road cracks.
The student is YOLOv11n-seg (2.8M parameters), the first teacher is SAM\,2.1 Hiera-Large (224M), and the data is Crack500.
The work has four steps:
\begin{enumerate}
\item Build and validate the method with SAM\,2. This is where I am now.
\item Choose which distillation losses to combine, and their weights.
\item Repeat with SAM\,1 and SAM\,3 as teachers.
\item Combine the teachers.
\end{enumerate}

\section{Step 1: SAM\,2 as teacher}
\subsection*{1a. What I have learned so far}
\begin{center}\small
\begin{tabular}{lc}
\toprule
Model (seed 0, 150 epochs) & Val mask mAP50\\
\midrule
Baseline, Ultralytics augmentation (A4) & \textbf{0.540}\\
Baseline, flips and rotations (A1) & 0.514\\
Gated KD from SAM\,2 logits (A5) & 0.511\\
Plain KD from SAM\,2 logits (A2) & 0.503\\
\bottomrule
\end{tabular}\end{center}
I added SAM's soft logits as an extra loss on images that already had ground-truth masks. Once the student was trained with
augmentation, this gave no gain. I measured why: on those images SAM without fine-tuning is less accurate than the labels
(median IoU 0.605), so it has nothing new to teach there.

\subsection*{1b. The method: SAM teaches on images without masks}
SAM is good at turning a cheap prompt, a box, into a mask. So I use the teacher on the images where the student has no mask.
\begin{enumerate}
\item Masks exist for only $p$\% of the train photos. The rest have boxes only. I take the boxes from the GT to simulate box annotation.
\item I fine-tune SAM's mask decoder on the $p$\% masks and keep the image encoder frozen.
\item SAM turns the boxes of the remaining crops into masks.
\item YOLOv11n-seg trains on the GT masks plus the SAM masks, with full augmentation. Validation and test always use GT.
\end{enumerate}
SAM's knowledge reaches the student through the masks it produces, which is called label distillation. These masks are ordinary labels,
so the student can train with mosaic and scale augmentation. Cached logits did not allow that.

\subsection*{1c. Experiments (Kaggle T4)}
\begin{center}\small
\begin{tabular}{llp{6.6cm}p{4.2cm}}
\toprule
Run & Masks & Training labels & Role\\
\midrule
R0 & 100\% & GT on 1896 crops (done: 0.540) & Upper bound\\
R1 & 25\% & GT on 478 crops only (done: 0.450) & Lower bound\\
R2 & 25\% & GT on 478 + SAM\,2 masks on 1418 (done: 0.538) & Main result\\
R3 & 0\% & SAM\,2 masks on all 1896, no fine-tuning (done: 0.335) & No masks at all\\
R4 & 25\% & GT on 478 + masks from the R1 student on 1418 & Same pipeline without a teacher\\
R5 & 10\% & as R1 and R2 with 201 labelled crops & Second point of the curve\\
\bottomrule
\end{tabular}\end{center}
R2 minus R1 is the gain from distillation. R2 minus R4 is the part of that gain that comes from SAM: R4 is the standard self-training
baseline, where the student labels the extra images alone, and reviewers usually ask for it.

I report val mask mAP50, test pixel F1, and the quality of the teacher's masks against the hidden GT.
I fixed the success criterion before running anything: R2 must beat R1 by at least 0.02 val mAP50, with the same sign on test pixel F1,
over 3 seeds.

\subsection*{1d. First results (seed 0)}
\begin{center}\small
\begin{tabular}{lcccc}
\toprule
Run & Val mAP50 & Test mAP50 & Test crop F1 & Photo F1\\
\midrule
R0, 100\% masks & 0.540 & 0.298 & 0.731 & 0.707\\
R1, 25\% masks & 0.450 & 0.257 & 0.724 & 0.712\\
R2, 25\% masks + SAM\,2 masks & 0.538 & 0.307 & 0.737 & 0.712\\
R3, SAM\,2 masks only, no fine-tuning & 0.335 & 0.125 & 0.444 & 0.425\\
\bottomrule
\end{tabular}\end{center}
With a quarter of the masks, SAM's labels bring the student back to the fully labelled model on mask mAP50: R2 gains 0.088 over R1 on
validation and 0.050 on the test crops. Fine-tuning the teacher is what makes this work. On the 478 labelled crops it raised SAM's
validation pixel F1 from 0.518 to 0.828, and its masks on the 1418 unlabelled crops reached F1 0.80 against the hidden GT.
Without fine-tuning SAM over-segments (precision 0.38), and R3 is far below R1.

The gain in pixel F1 is small. R2 is 0.013 above R1 on the test crops and equal to it on the full photos. Pixel F1 hardly depends on the
number of masks: R1 is already within 0.01 of R0, so there was little to recover on that metric.

These are single-seed numbers. R4 is running, and R5 and the seeds for R1 and R2 come next. I chose the teacher's best epoch on the same
validation masks that score the student; its F1 stayed between 0.81 and 0.83 over all ten epochs, so I expect the effect to be small.
The first R2 run stopped before training because the download of the SAM checkpoint from Meta's server failed; the notebooks now fall back
to the Hugging Face copy.

\section{Step 2: which losses to combine, and their weights}
The code has four terms: the task loss (YOLO box, class and mask), mask KD (KL to SAM logits), feature KD (channel-wise, on neck levels P3 to P5)
and a boundary loss. I tuned the current weights (0.96, 1.87, 0.81) before the code audit, so I have to choose them again.
Starting from the R2 model:
\begin{enumerate}
\item Add one term at a time: mask KD, then boundary, then feature KD (3 runs).
\item Set each weight by scale, so the term is about 25\% of the task loss in the first epochs. For the best term, also try half and double that weight (2 runs).
\item Combine the terms that gained at least 0.01 alone (1 run), then run 3 seeds.
\end{enumerate}
I make every choice on validation and use the test set once at the end. One technical constraint: cached logits and features cannot follow
mosaic, so I apply these terms either with flip and rotation augmentation or in the final epochs, where mosaic is off.

\section{Step 3: other teachers, SAM\,1 and SAM\,3}
The pipeline of Step 1 stays the same and only the teacher changes. For each teacher I report its own mask quality on validation
and the score of the student trained from it, in the R2 setting.
\begin{center}\small
\begin{tabular}{lp{5.2cm}p{6.6cm}}
\toprule
Teacher & Prompt & Note\\
\midrule
SAM\,2.1 Hiera-L & box & Step 1\\
SAM\,1 ViT-H & box & Older, heavier image encoder; same box-prompt interface\\
SAM\,3 & box, and the text ``crack'' & The text prompt needs no boxes, so it could label images that have no annotation at all\\
\bottomrule
\end{tabular}\end{center}
\sloppy One paper reports about 56\% crack IoU for SAM\,3 with the prompt ``crack'' and no fine-tuning
(\href{https://arxiv.org/abs/2607.12292}{arXiv:2607.12292}), so I expect fine-tuning to matter for every teacher.

\section{Step 4: combining the teachers}
The teachers make different errors, so I can fuse their masks before training the student. I will try three rules:
\begin{itemize}
\item Average the teachers' probabilities per pixel, after putting their logits on a common scale.
\item Keep a crack pixel only if at least two teachers predict it, which removes false positives.
\item Per pixel, trust the more confident teacher.
\end{itemize}
I evaluate the fused teacher like a single one (its mask quality, then the student trained from it) and compare it with the best single teacher.

\section{Expected outcome}
The result is a distillation method for a 2.8M-parameter real-time model. It specifies how SAM's knowledge is transferred
(masks on unlabelled images, plus the loss terms of Step 2), compares SAM\,1, SAM\,2 and SAM\,3 as teachers, and gives a rule for fusing them.
The study also reports the earlier finding that logit KD on labelled images gives no gain under matched augmentation.

\paragraph{Limitations so far.} Every result is from one seed. I derive the boxes from the GT masks, so they are tighter than hand-drawn boxes;
the code has a box-jitter option to test this. Steps 3 and 4 are not implemented yet.

\sloppy\paragraph{Related work.} I did a short search and read abstracts only.
Label distillation from foundation models into small students appears in \href{https://arxiv.org/abs/2505.21904}{CAST (arXiv:2505.21904)},
\href{https://arxiv.org/abs/2604.03841}{arXiv:2604.03841} and \href{https://arxiv.org/abs/2411.16196}{arXiv:2411.16196}.
\href{https://arxiv.org/abs/2407.13553}{arXiv:2407.13553} uses SAM box prompts to make pseudo-labels in medical imaging.
I found no crack paper that distils SAM box-prompted masks into a real-time model, and none that compares several SAM versions as teachers.
\end{document}
